\begin{theorem}[Semi-regularity and the group algebra]
    \label{thm:peps_semiregular_group_algebra}
    \lean{Representation.linearIndependent_iff_injective_asAlgebraHom,
        Representation.isSemiRegular_of_linearIndependent,
        Representation.isSemiRegular_iff_linearIndependent,
        Representation.isSemiRegular_iff_injective_asAlgebraHom}
    \leanok
    \uses{def:peps_semiregular}
    Let $G$ be a finite group and let $\rho$ be a finite-dimensional complex
    representation. The following assertions are equivalent:
    $\rho$ contains every irreducible representation; the operators
    $(\rho(g))_{g\in G}$ are linearly independent; and the induced algebra
    homomorphism $\mathbb C[G]\to\operatorname{End}(V)$ is injective.
    This is an algebraic characterization of the occurrence condition in
    \cite[Definition~4.5 and Lemma~4.6]{Schuch2010PEPS}.
    More generally, for a representation of a monoid over a commutative
    semiring, independence of its operators is equivalent to injectivity
    of its monoid-algebra homomorphism.
\end{theorem}
\begin{proof}\leanok
    \uses{thm:peps_semiregular_trace}
    The trace-dual identity gives independence for a semi-regular
    representation. Conversely, if an irreducible character $\chi$ is
    absent, its character projector vanishes:
    \begin{align}
        \frac{\chi(1)}{|G|}\sum_{g\in G}\chi(g^{-1})\rho(g)=0.
        \label{eq:peps_absent_character}
    \end{align}
    Indeed, the projector vanishes on every irreducible summand of $V$.
    Independence in (\ref{eq:peps_absent_character}) would imply
    $\chi(1)=0$, contrary to the positive dimension of the irreducible
    representation. The standard basis of the group algebra identifies
    the last assertion with independence of the group operators.
\end{proof}

\begin{theorem}[Reversing a semi-regular bond]
    \label{thm:peps_semiregular_dual}
    \lean{Representation.IsSemiRegular.dual,
        Representation.isSemiRegular_dual_iff}
    \leanok
    \uses{def:peps_semiregular}
    A finite-dimensional complex representation $\rho$ of a finite group
    is semi-regular if and only if its contragredient
    $\rho^*(g)=\rho(g^{-1})^{\mathsf T}$ is semi-regular.
    This justifies reversing the bond orientation in the contraction
    convention of \cite[Definition~5.1 and Lemma~5.2]{Schuch2010PEPS}.
\end{theorem}
\begin{proof}\leanok
    \uses{thm:peps_semiregular_group_algebra}
    Inversion permutes the group indices, and transposition is injective
    on endomorphisms of a finite-dimensional vector space. It therefore
    preserves independence of the group operators. Conversely,
    independence of their transposes implies independence of the
    original operators.
\end{proof}

\begin{theorem}[Nonzero tensor factors of a semi-regular representation]
    \label{thm:peps_semiregular_nonzero_tensor}
    \lean{Representation.IsSemiRegular.tprod_of_nontrivial,
        Representation.IsSemiRegular.tprod_left_of_nontrivial}
    \leanok
    \uses{def:peps_semiregular}
    Let $G$ be a finite group, let $\rho$ be a finite-dimensional
    semi-regular complex representation, and let $\sigma$ be any
    nonzero finite-dimensional complex representation. Then both
    $\rho\otimes\sigma$ and $\sigma\otimes\rho$ are semi-regular.
    This strengthens the tensor-product assertion in
    \cite[Lemma~5.2]{Schuch2010PEPS}: the second factor need not contain
    the trivial representation.
\end{theorem}
\begin{proof}\leanok
    \uses{thm:peps_semiregular_group_algebra, thm:peps_semiregular_trace,
        thm:peps_semiregular_equivalence}
    A dependence among the tensor-product operators is a dependence
    among the tensors $\rho(g)\otimes\sigma(g)$, by the natural
    identification of their endomorphism spaces. Applying the
    trace-dual functional at $h$ to the first factor gives
    $c_h\sigma(h)=0$. Since $\sigma(h)$ is invertible on a nonzero
    vector space, $c_h=0$. The group-algebra criterion proves
    semi-regularity. Interchanging the tensor factors gives the
    other order.
\end{proof}
